\chapter{The Middle Perverse \texorpdfstring{\(t\)}{t}-Structure}
\label{chap:perverse-t-structure}

Ordinary cohomological degree measures how far a complex is from being a
sheaf.  Perverse degree measures that failure after compensating for the
dimension of the support.  This compensation is exactly what makes a local
system on a smooth \(d\)-fold live in degree \(-d\), makes Verdier duality
\(t\)-exact, and makes intersection cohomology an object of an abelian
category.

Throughout this chapter \(X\) is a finite-type scheme over a field, and
\(\Dbc(X)=\Dbc(X,\Qbarl)\) is the bounded constructible derived category.  The
comparison with the classical constructible adic category is compatible with
stalks, costalks, and recollement, so this is the usual rational middle
perversity, now realized on \(X_{\proet}\).  Integral coefficients require the
distinct \(p\) and \(p+\) structures of \cref{chap:integral-perverse}; we make
no integral self-duality claim in the present chapter.  For a geometric point
\(\bar x\to X\), let
\[
  \delta(x)=\dim\overline{\{x\}}.
\]

\section{Support and cosupport conditions}

\begin{definition}[Middle perversity]
\label{def:middle-perversity}
Define full subcategories of \(\Dbc(X)\) by
\begin{align*}
 {}^p\D^{\leq0}(X)
 &=\{\cK:H^a(i_{\bar x}^*\cK)=0
       \text{ whenever }a>-\delta(x)\},\\
 {}^p\D^{\geq0}(X)
 &=\{\cK:H^a(i_{\bar x}^!\cK)=0
       \text{ whenever }a<-\delta(x)\}.
\end{align*}
Here \(i_{\bar x}:\bar x\to X\) is the geometric-point morphism.  Put
\[
 {}^p\D^{\leq n}={}^p\D^{\leq0}[-n],\qquad
 {}^p\D^{\geq n}={}^p\D^{\geq0}[-n].
\]
\end{definition}

The stalk condition has an equivalent global form:
\[
 \cK\in{}^p\D^{\leq0}
 \quad\Longleftrightarrow\quad
 \dim\Supp\cH^a(\cK)\leq-a\quad\text{for all }a.
\tag{8.1}\label{eq:perverse-support}
\]
The costalk condition is not obtained by replacing stalks with ordinary
stalks.  It records local cohomology transverse to the strata, and is the
Verdier-dual condition.

\begin{theorem}[BBD middle \(t\)-structure]
\label{thm:middle-t-structure}
The pair
\[
  ({}^p\D^{\leq0}(X),{}^p\D^{\geq0}(X))
\]
is a bounded \(t\)-structure on \(\Dbc(X)\).  Its heart
\[
  \Perv(X,\Qbarl)
  ={}^p\D^{\leq0}(X)\cap{}^p\D^{\geq0}(X)
\]
is an abelian category.  The truncation functors are denoted
\(\ptau_{\leq n},\ptau_{\geq n}\), and
\[
  \pH^n(\cK)=\bigl(\ptau_{\geq n}\ptau_{\leq n}\cK\bigr)[n]
\]
is the \(n\)-th perverse cohomology object.
\end{theorem}

\begin{proof}
Fix first a finite algebraic stratification
\(\mathscr S=\{S_\alpha\}\) of \(X\) by smooth, locally closed,
pure-dimensional strata, and let \(\D^b_{\mathscr S}(X)\) be the full
subcategory of complexes whose cohomology sheaves are constructible along
\(\mathscr S\).  We construct the \(t\)-structure on this fixed category by
induction on the number of strata.

If there is one stratum of dimension \(d\), all cohomology sheaves are lisse
and the required \(t\)-structure is the ordinary lisse \(t\)-structure shifted
by \(d\).  In general let \(j:U\hookrightarrow X\) be the union of the open
strata and let \(i:Z\hookrightarrow X\) be its closed complement.  On each
connected component of \(U\) use the ordinary lisse \(t\)-structure shifted by
the dimension of that component; on \(Z\) use the structure already obtained
by induction.  Recollement glues them: the two halves are characterized by
\[
 \begin{split}
 K\in{}^p\D_{\mathscr S}^{\leq0}(X)
 &\Longleftrightarrow
 j^*K\in{}^p\D^{\leq0}(U),\quad i^*K\in{}^p\D^{\leq0}(Z),\\
 K\in{}^p\D_{\mathscr S}^{\geq0}(X)
 &\Longleftrightarrow
 j^*K\in{}^p\D^{\geq0}(U),\quad i^!K\in{}^p\D^{\geq0}(Z).
 \end{split}
\]
The two recollement triangles construct the truncation triangle, while the
adjunctions give
\[
 \Hom({}^p\D^{\leq0}_{\mathscr S},
      {}^p\D^{\geq1}_{\mathscr S})=0.
\]
Induction on the strata identifies these glued conditions with the stated
pointwise stalk and costalk inequalities.

If \(\mathscr S'\) refines \(\mathscr S\), the two constructions agree on
\(\D^b_{\mathscr S}(X)\), because the pointwise inequalities do not mention a
stratification.  Every finite family of constructible complexes is adapted to
a common finite refinement.  We may therefore pass from the fixed-stratum
categories to \(\Dbc(X)\); the orthogonality and truncation axioms are checked
inside such a common category.  Boundedness follows from bounded
constructibility and finite dimension.  Finally, the heart of any bounded
\(t\)-structure is abelian, with kernels and cokernels obtained from the
perverse cohomology of a cone.  This is the gluing construction of
BBD, Sections 1.4 and 2.1--2.2, for the dimension function \(-\delta\);
see \cite{BBD}.
\end{proof}

\begin{perversebox}
The proof uses only constructibility, recollement, and the local dimension
function.  It contains no weights and therefore lies safely upstream of the
Weil conjectures.
\end{perversebox}

\section{The dimension shift}

\begin{proposition}[Local systems are perverse after the dimension shift]
\label{prop:lisse-shift-perverse}
Let \(X\) be smooth and pure of dimension \(d\), and let \(\cL\) be lisse.
Then \(\cL[d]\) is perverse.
\end{proposition}

\begin{proof}
For a point \(x\) of dimension \(e\), the stalk of \(\cL[d]\) is concentrated
in degree \(-d\leq-e\), which gives the support inequality.  Purity of the
smooth local ring and absolute cohomological purity identify the transverse
costalk, up to a lisse factor, with a Tate twist placed in the codimension
predicted by \(d-e\).  Consequently it has no cohomology below \(-e\), giving
the cosupport inequality.  Equivalently, choose a smooth stratification and
verify the two inequalities on normal slices.  Both are equalities on the
generic stratum.
\end{proof}

\begin{example}
On a smooth curve \(C\), a lisse sheaf \(\cL[1]\) is perverse.  At a closed
point its stalk is concentrated in degree \(-1\), and its costalk in degree
\(1\).  A general perverse sheaf on a curve may have stalk cohomology in
degrees \(-1,0\) and costalk cohomology in degrees \(0,1\); the degree-zero
boundary terms detect punctual subobjects and quotients.  This is precisely
the information controlled by intermediate extension.
\end{example}

\begin{example}
For a closed point \(i:x\hookrightarrow X\), the skyscraper \(i_*V\) is
perverse with no shift.  More generally a finite map sends perverse sheaves to
perverse sheaves.  Thus objects supported in different dimensions naturally
occupy different ordinary cohomological degrees but the same perverse degree.
\end{example}

\begin{warning}
If \(X\) is singular, the shifted constant sheaf \(\Qbarl[\dim X]\) need not be
perverse.  The canonical replacement is the intersection complex constructed
in \cref{chap:intermediate-extension}.  The equality
\(\IC_X=\Qbarl[\dim X]\) is a smoothness statement, not a definition.
\end{warning}

\section{Recollement in the heart}

Let \(i:Z\hookrightarrow X\) be closed and
\(j:U=X\setminus Z\hookrightarrow X\) its open complement.  The functors
\(j^*=j^!\) and \(i_*=i_!\) are perverse \(t\)-exact.  The four remaining
functors are one-sided exact:
\[
 \begin{array}{c|c}
 \text{functor}&\text{perverse exactness}\\
 \hline
 i^*,\ j_!&\text{right \(t\)-exact}\\
 i^!,\ Rj_*&\text{left \(t\)-exact}.
 \end{array}
\tag{8.2}\label{eq:recollement-exactness}
\]
Here “right \(t\)-exact” means preserving
\({}^p\D^{\leq0}\); “left \(t\)-exact” means preserving
\({}^p\D^{\geq0}\).  This convention will always be stated when used.

Define heart-level functors
\[
 {}^pj_!=\pH^0j_!,\qquad
 {}^pj_*=\pH^0Rj_*,\qquad
 {}^pi^*=\pH^0i^*,\qquad
 {}^pi^!=\pH^0i^!.
\]
For \(\cP\in\Perv(X)\), the recollement triangles yield morphisms in the
perverse heart
\begin{align}
 {}^pj_!j^*\cP&\longrightarrow\cP\longrightarrow
      i_*{}^pi^*\cP\longrightarrow0,
\tag{8.3}\label{eq:perverse-recollement-one}\\
 0&\longrightarrow i_*{}^pi^!\cP\longrightarrow\cP
      \longrightarrow{}^pj_*j^*\cP.
\tag{8.4}\label{eq:perverse-recollement-two}
\end{align}
Together with their neighboring perverse cohomology terms, these are exact
sequences.  They detect subobjects and quotients supported on the boundary.

\section{Verdier duality}

\begin{theorem}[Self-duality of middle perversity]
\label{thm:perverse-verdier-duality}
Verdier duality exchanges the two halves of the middle \(t\)-structure:
\[
 \D_X({}^p\D^{\leq0})={}^p\D^{\geq0},\qquad
 \D_X({}^p\D^{\geq0})={}^p\D^{\leq0}.
\]
It therefore restricts to an exact contravariant equivalence
\[
 \D_X:\Perv(X)^{\op}\xrightarrow{\sim}\Perv(X)
\]
and satisfies
\[
 \D_X\pH^a(\cK)\simeq\pH^{-a}(\D_X\cK).
\tag{8.5}\label{eq:duality-perverse-cohomology}
\]
\end{theorem}

\begin{proof}
Local Verdier duality exchanges a stalk with the dual of the corresponding
costalk.  The inequality \(a>-\delta(x)\) is therefore changed into
\(-a<-\delta(x)\), precisely the other defining condition.  Biduality gives
equality rather than mere inclusion.  General \(t\)-structure formalism then
gives \eqref{eq:duality-perverse-cohomology}.
\end{proof}

This self-duality is the reason for using the middle perversity.  Ordinary
sheaves are not stable under Verdier duality on a positive-dimensional space.

\section{Functorial amplitudes}

The following list is used repeatedly.  It is best viewed as a table of
interfaces, not as permission to suppress shifts.

\begin{proposition}[Basic perverse exactness]
\label{prop:basic-perverse-exactness}
Let all morphisms be separated and of finite type.
\begin{enumerate}
\item If \(f:X\to Y\) is smooth of pure relative dimension \(r\), then
      \(f^*[r]\) and \(f^![-r]=f^*[r](r)\) are \(t\)-exact.
\item If \(f\) is finite, then \(Rf_*=Rf_!\) is \(t\)-exact.
\item If \(f\) is affine, then \(Rf_*\) is right \(t\)-exact and \(Rf_!\) is
      left \(t\)-exact; equivalently, for \(\cP\in\Perv(X)\),
      \[
        \pH^a(Rf_*\cP)=0\ (a>0),\qquad
        \pH^a(Rf_!\cP)=0\ (a<0).
      \]
      If \(f\) is also quasi-finite, both functors are \(t\)-exact.
\item External product is \(t\)-exact on perverse hearts:
      \(\cP\boxtimes\cQ\in\Perv(X\times Y)\).
\end{enumerate}
\end{proposition}

\begin{proof}
The smooth statement follows by comparing dimensions of strata and using
smooth base change for stalks and costalks.  For a finite map, dimensions of
closures are preserved and proper base change computes both conditions
fiberwise.  The affine statement is the perverse form of Artin vanishing.  One
embeds an affine morphism into affine space over the base, filters by generic
hyperplane sections, and applies ordinary Artin vanishing to the successive
normal slices; Verdier duality gives the compact-support half.  For an affine
quasi-finite morphism the relative support dimension is zero, so the two
one-sided bounds meet.  External products follow from the Kunneth formula for
stalks and costalks and addition of dimensions.  The affine and quasi-finite
assertions are BBD 4.1.1--4.1.3; smooth pullback is BBD 4.2.4, and external
products follow from BBD 4.2.7--4.2.8 (or directly from the stalk/costalk
Kunneth formula).
\end{proof}

\begin{corollary}[Artin vanishing for perverse sheaves]
\label{cor:perverse-artin-vanishing}
If \(X\) is affine over an algebraically closed field and
\(\cP\in\Perv(X)\), then
\[
 H^a(X,\cP)=0\quad(a>0),\qquad
 H_c^a(X,\cP)=0\quad(a<0).
\tag{8.6}\label{eq:perverse-artin}
\]
\end{corollary}

\begin{proof}
Apply \cref{prop:basic-perverse-exactness} to the affine structural morphism
\(X\to\Spec k\).  On a point the perverse and ordinary \(t\)-structures agree.
\end{proof}

\section{Perverse spectral sequences without degeneration}

For any cohomological functor \(T:\Dbc(X)\to\mathcal A\), the perverse
truncation tower gives a spectral sequence
\[
 E_2^{a,b}=T^{a}(\pH^b\cK)
 \Longrightarrow T^{a+b}(\cK).
\tag{8.7}\label{eq:perverse-spectral-sequence}
\]
Because \(\cK\) is bounded, the induced filtration is finite.  Every graded
piece of the abutment is a subquotient of an \(E_2\)-term.  This is all that is
needed for weight inequalities: upper and lower weight bounds pass to
subquotients.  No degeneration assertion is used in the direct-image theorem.

\begin{formalization}
The reusable formalization layer should expose:
\begin{enumerate}
\item a dimension function on points and the two pointwise predicates in
      \cref{def:middle-perversity};
\item the glued bounded \(t\)-structure and its heart;
\item perverse truncation triangles and functorial perverse cohomology;
\item exact recollement functors and the one-sided amplitude table;
\item a natural isomorphism expressing Verdier self-duality;
\item the finite filtration attached to
      \eqref{eq:perverse-spectral-sequence}, separated from any degeneration
      theorem.
\end{enumerate}
The point case and the smooth \(d\)-fold shift should be executable tests:
\(\Perv(\Spec k)\) is the ordinary sheaf category and \(\cL[d]\) is perverse.
\end{formalization}
